\section{Exceptional terms and the power-saving induction}
\label{sec:induction}

We now combine the shorter-scale comparison with cancellation on good
residue assignments to obtain a recurrence for $\mathcal Y_l(N)$.
The exceptional parts of the prime laws contribute a summable family of
additional shorter-scale terms.  The expansion and contraction argument
adapt \cite[Proposition~9.2 and Theorem~9.3]{squarepower}.

Retain all parameters from \eqref{eq:14}.  For $B\subseteq\mathcal P_N$, put
\[
 D_B=M_0q_B,\qquad
 n_B=\left\lceil\frac{N}{\lambda(D_B)}\right\rceil,
 \qquad
 \eta_{l,B}=\prod_{p\in B}\eta_{l,p}.
\]
Thus $q_\varnothing=\eta_{l,\varnothing}=1$ and
$D_\varnothing=M_0$.  Recall that $\rho_l$ is the proportion of good
residue assignments for $h_l$, and that
$\rho_l\ge\rho_*=M_0^{1-L}>0$ by
Lemma~\ref{transfer:good-assignments}.

\begin{proposition}[Recurrence for the maximal functional]
\label{induction:recurrence}
There are constants $C_1\ge1$ and $N_1\ge3$, depending only on the fixed
polynomial $h$, its chosen roots, and the fixed structural choices, such that
for every integer $N\ge N_1$ and every positive integer $l\le N^\zeta$,
\begin{equation}
 \begin{split}
 \mathcal Y_l(N)\le C_1N^{-\sigma}
 +(1+u_N)\Bigg(&
   (1-\rho_l)\mathcal Y_{lD_\varnothing}(n_\varnothing)\\
 &+\rho_l
   \sum_{\substack{\varnothing\ne B\subseteq\mathcal P_N\\q_B\le N^\tau}}
      \eta_{l,B}\mathcal Y_{lD_B}(n_B)\Bigg),
 \end{split}
 \label{eq:31}
\end{equation}
where
\[
 u_N=\bigl(1+M_0^kN^{k\tau-1}\bigr)^d-1\longrightarrow0.
\]
Every scale $n_B$ appearing on the right belongs to $[1,N-1]$.
\end{proposition}

\begin{proof}
Fix $A\subseteq[N]$ satisfying \eqref{eq:24}.  Use the functions $f_s$ and
their lifts $G_s$ from Section~\ref{sec:transfer}, and, for a residue
assignment $\mathbf s=(s_v)_{v\in V}$, write
\[
 \Phi_{\mathbf s}((z_v)_{v\in V})=\prod_{v\in V}G_{s_v}(z_v).
\]
Since $\E_s f_s=\1_A$, linearity of the lift and multilinearity of the
integral give
\begin{equation}
 Y_l(N,A)=\E_{\mathbf s}
   \int\Phi_{\mathbf s}\,d\bigotimes_{p\in\mathcal P_N}\mu_{l,p},
 \label{induction:residue-average}
\end{equation}
where the residues $s_v$ are independent and uniform modulo $M_0$.

For each $B\subseteq\mathcal P_N$, define the nonnegative measure
\[
 \kappa_{l,B}=
   \bigotimes_{p\in B}\eps_{l,p}
   \otimes\bigotimes_{p\in\mathcal P_N\setminus B}\mu_{l,p}.
\]
Its mass is $\eta_{l,B}$.  Expanding the finite product gives the identity
of signed measures
\begin{equation}
 \bigotimes_{p\in\mathcal P_N}\mu_{l,p}
 -\bigotimes_{p\in\mathcal P_N}(\mu_{l,p}-\eps_{l,p})
 =\sum_{\varnothing\ne B\subseteq\mathcal P_N}
      (-1)^{|B|+1}\kappa_{l,B}.
 \label{induction:exceptional-expansion}
\end{equation}
The complement of $B$ therefore carries the full laws $\mu_{l,p}$, to
which Proposition~\ref{transfer:shorter-scale} applies.  The functions
$\Phi_{\mathbf s}$ can have either sign; we will take the absolute value
of each integral on the right of
\eqref{induction:exceptional-expansion}.

The exceptional masses have a uniform bound with an extra factor $q_B$:
\begin{equation}
 \sum_{\varnothing\ne B\subseteq\mathcal P_N}q_B\eta_{l,B}
 =\prod_{p\in\mathcal P_N}(1+p\eta_{l,p})-1
 \le\prod_{p\ge P}(1+p^{-3})-1\le\frac14.
 \label{eq:32}
\end{equation}
Here we used $\eta_{l,p}\le p^{-4}$ and the choice of $P$ in
\eqref{eq:9}.  This bound first allows us to discard the terms with
$q_B>N^\tau$.

To see this, suppose $\eta_{l,B}>0$.  Simultaneous translation invariance
of $\eps_{l,p}$ implies that each single-label marginal of
$\eps_{l,p}/\eta_{l,p}$ is uniform on $\F_p$.  The same holds for the full
laws.  Thus each single-slot marginal of $\kappa_{l,B}/\eta_{l,B}$ is
uniform on $X_{\mathcal P_N}$.  Ordinary H\"older's inequality and the
moment estimate \eqref{eq:12}, at the fixed even exponent $d$, give
\begin{equation}
 \left|\int\Phi_{\mathbf s}\,d\kappa_{l,B}\right|
 \le\eta_{l,B}\prod_{v\in V}\|G_{s_v}\|_d
 \le C_\epsilon\eta_{l,B}N^\epsilon
 \qquad(\epsilon>0).
 \label{induction:exceptional-moment}
\end{equation}
The constants are uniform in $l,B,A$ and $\mathbf s$; the bound
$|f_s|\le M_0$ only introduces a fixed factor.  Terms of zero mass vanish.
Combining \eqref{induction:exceptional-moment} with \eqref{eq:32} yields
\[
 \sum_{\substack{\varnothing\ne B\subseteq\mathcal P_N\\q_B>N^\tau}}
   \left|\int\Phi_{\mathbf s}\,d\kappa_{l,B}\right|
 \le \frac{C_\epsilon}{4}N^{-\tau+\epsilon}
 =O(N^{-\sigma}),
\]
where, for example, $\epsilon=\tau/2$ suffices because
$\sigma=\gamma\tau/4<\tau/2$.

For $q_B\le N^\tau$ and positive mass, condition instead on all the
labels at primes in $B$.  The laws at the remaining primes are unchanged,
and the function in each slot is the corresponding specialization of
$G_{s_v}$.  Proposition~\ref{transfer:shorter-scale}, which allows arbitrary
slotwise specializations, gives
\begin{equation}
 \left|\int\Phi_{\mathbf s}\,d\kappa_{l,B}\right|
 \le\eta_{l,B}\left(
      \Lambda_B^d\mathcal Y_{lD_B}(n_B)+C_2N^{-\sigma}\right),
 \qquad
 \Lambda_B=\frac{\lambda(D_B)n_B}{N}.
 \label{induction:exceptional-transfer}
\end{equation}
The bound is uniform in every conditioned tuple, so integration over
the prescribed labels preserves it.  It also holds for zero-mass terms.
By \eqref{eq:27}, $\Lambda_B^d\le1+u_N$; by \eqref{eq:32},
$\sum_{B\ne\varnothing}\eta_{l,B}\le1/4$.  Therefore summing the errors
in \eqref{induction:exceptional-transfer} still costs $O(N^{-\sigma})$.

For a good assignment $\mathbf s$, Proposition~\ref{transfer:cancellation}
bounds the integral against
$\bigotimes_p(\mu_{l,p}-\eps_{l,p})$ by $O(N^{-\sigma})$ in absolute
value.  The expansion \eqref{induction:exceptional-expansion} and the
two exceptional estimates consequently give
\[
 \int\Phi_{\mathbf s}\,d\bigotimes_p\mu_{l,p}
 \le O(N^{-\sigma})+(1+u_N)
   \sum_{\substack{\varnothing\ne B\subseteq\mathcal P_N\\q_B\le N^\tau}}
     \eta_{l,B}\mathcal Y_{lD_B}(n_B).
\]
For every other assignment, apply
Proposition~\ref{transfer:shorter-scale} directly with $B=\varnothing$.
The resulting bound is
$(1+u_N)\mathcal Y_{lM_0}(n_\varnothing)+O(N^{-\sigma})$.
Averaging these estimates in \eqref{induction:residue-average} gives
the coefficients $\rho_l$ and $1-\rho_l$ in \eqref{eq:31}.
The bounds are uniform in $A$, so we may take the maximum defining
$\mathcal Y_l(N)$.

Finally, $k\tau<1$ by \eqref{eq:14}, so $u_N\to0$, and the bound
$1\le n_B<N$ is part of Proposition~\ref{transfer:shorter-scale}.
\end{proof}

The recurrence decreases $N$ while replacing $l$ by $lD_B$.  Its
cancellation hypothesis $l\le N^\zeta$ need not persist at the smaller
scale.  The next estimate is therefore stated for every $l$: a small
positive power of $l$ permits the uniform moment bound to handle the range
where that hypothesis fails.

\begin{proposition}[Decay throughout the auxiliary family]
\label{induction:decay}
There are constants $a_0,b_0>0$ and $C_0\ge1$, depending only on the fixed
polynomial $h$, its chosen roots, and the fixed structural choices, such that
for all positive integers $l,N$,
\begin{equation}
 \mathcal Y_l(N)\le C_0l^{b_0}N^{-a_0}.
 \label{eq:33}
\end{equation}
The exponents $a_0,b_0$ can be chosen as functions only of the structural
parameters $k,\zeta,\sigma,M_0,\rho_*$.  If the chosen roots of $h$
satisfy Condition~\eqref{local:unit-condition} and the choices depending
only on $k$ from Sections~\ref{sec:local} and~\ref{sec:tuples} are used, then
$a_0,b_0,d$ depend only on $k$, while $C_0$ may still depend on the fixed data.
\end{proposition}

\begin{proof}
To select the exponents, suppose that at a scale where the recurrence applies
the trial bound
$\mathcal Y_{lD_B}(n_B)\le C(lD_B)^b n_B^{-a}$ is available at every
successor scale, with $C\ge1$ and $a,b>0$.  Put $s=ka+b$.  Since
$n_B\ge N/\lambda(D_B)$ and $\lambda(D_B)\le D_B^k$, each trial bound gives
\[
 C(lD_B)^b n_B^{-a}
 \le C l^bN^{-a}\lambda(D_B)^aD_B^b
 \le C l^bN^{-a}D_B^s.
\]
If $s\le1$, then $q_B^s\le q_B$, so Equation~\eqref{eq:32} bounds the
coefficient of $C l^bN^{-a}$ inside the
parentheses of \eqref{eq:31} by
\[
 M_0^s\left(1-\rho_l+
   \rho_l\sum_{\substack{\varnothing\ne B\subseteq\mathcal P_N\\q_B\le N^\tau}}
      \eta_{l,B}q_B^s\right)
 \le M_0^s(1-3\rho_*/4).
\]
This upper bound tends to $1-3\rho_*/4<1$ as $s$ tends to zero.  Two
further restrictions handle the other parts of the argument.
For any proposed $a>0$, Corollary~\ref{lift:density} gives a bound
$\mathcal Y_l(N)\le C_aN^a$ for all positive $l,N$.  Taking
$b=2a/\zeta$ makes $l^bN^{-a}>N^a$ whenever $l>N^\zeta$, so this
subpower bound covers the range outside the recurrence after increasing
$C$.  Finally, $a<\sigma$ makes $C_1N^{-\sigma}=o(N^{-a})$, allowing
the additive error to be absorbed at sufficiently large scales.

Choose explicitly
\[
 a_0=\frac12\min\left\{
   \sigma,\quad \frac1{k+2/\zeta},\quad
   \frac{-\log(1-3\rho_*/4)}{(k+2/\zeta)\log M_0}
 \right\},
 \qquad b_0=\frac{2a_0}{\zeta}.
\]
These numbers are positive because $M_0\ge2$ and $0<\rho_*\le1$.
They satisfy
\begin{equation}
 a_0<\sigma,\qquad
 s_0:=ka_0+b_0\le1,\qquad
 \theta:=M_0^{s_0}(1-3\rho_*/4)<1.
 \label{induction:exponent-choice}
\end{equation}
Indeed, $s_0=(k+2/\zeta)a_0\le1/2$, and the third term in the minimum gives
\[
 s_0\log M_0\le-\tfrac12\log(1-3\rho_*/4),
 \qquad
 \theta\le(1-3\rho_*/4)^{1/2}<1.
\]
This displays the asserted structural dependence.  Under
Condition~\eqref{local:unit-condition}, Sections~\ref{sec:local}
and~\ref{sec:tuples} permit $L,P,M_0,d$ to be fixed in terms of $k$ only.
Then $\rho_*=M_0^{1-L}$ and all parameters in \eqref{eq:14} depend only
on $k$, so the displayed choices of $a_0,b_0$ do as well.  The constants
in the estimates below, and consequently $C_0$, may depend on the fixed data.
All quantities in \eqref{induction:exponent-choice} are fixed from now on.

Corollary~\ref{lift:density} supplies a constant
$C_*\ge1$ such that
\begin{equation}
 \mathcal Y_l(N)\le C_*N^{a_0}
 \qquad(l,N\ge1).
 \label{induction:subpower-bound}
\end{equation}
If $l>N^\zeta$, then
\[
 l^{b_0}N^{-a_0}>N^{\zeta b_0-a_0}=N^{a_0},
\]
so \eqref{eq:33} follows from \eqref{induction:subpower-bound} for any
$C_0\ge C_*$.  This establishes the estimate outside the range of
the recurrence at every scale.

We choose the initial range before fixing $C_0$.  Set
$\kappa=(1-\theta)/2>0$, and take an integer $N_*\ge N_1$ so large that
for every $N\ge N_*$,
\begin{equation}
 \theta u_N\le\kappa,\qquad
 C_1N^{a_0-\sigma}\le\kappa.
 \label{induction:threshold}
\end{equation}
This is possible because $u_N\to0$ and $a_0<\sigma$; the choice is
independent of $C_0$.  Now choose
$C_0\ge\max\{1,C_*N_*^{2a_0}\}$.  For all $l\ge1$ and $N<N_*$,
\eqref{induction:subpower-bound} gives
\[
 \mathcal Y_l(N)\le C_*N^{a_0}
 \le C_0N^{-a_0}\le C_0l^{b_0}N^{-a_0}.
\]
Thus the initial range is covered simultaneously for every $l$.

Proceed by strong induction on $N$, simultaneously for all positive
integers $l$.  Suppose $N\ge N_*$ and the estimate is known at every
smaller scale.  The case $l>N^\zeta$ was established above, so assume
$l\le N^\zeta$ and apply Proposition~\ref{induction:recurrence}.
Its successor scales satisfy $1\le n_B<N$, so the simultaneous induction
hypothesis gives
$\mathcal Y_{lD_B}(n_B)\le C_0(lD_B)^{b_0}n_B^{-a_0}$ at every positive
index $lD_B$, with no restriction on $lD_B$ relative to $n_B$.
The trial calculation above, with $C=C_0$, $a=a_0$, $b=b_0$, and $s=s_0$,
therefore bounds the successor expression inside the parentheses of
\eqref{eq:31} by $\theta C_0l^{b_0}N^{-a_0}$.
Substituting in \eqref{eq:31}, and then using
\eqref{induction:threshold}, yields
\begin{align*}
 \mathcal Y_l(N)
 &\le C_1N^{-\sigma}
       +(1+u_N)\theta C_0l^{b_0}N^{-a_0}\\
 &\le \kappa N^{-a_0}
       +(1-\kappa)C_0l^{b_0}N^{-a_0}
 \le C_0l^{b_0}N^{-a_0}.
\end{align*}
The last step uses $C_0l^{b_0}\ge1$.  This closes the induction.
\end{proof}

\begin{proof}[Proof of Theorem~\ref{thm:main}]
Fix $k\ge2$.  For this degree, fix the numerical choices furnished under
Condition~\eqref{local:unit-condition} by Sections~\ref{sec:local},
\ref{sec:tuples}, and~\ref{sec:lifts}.  In particular,
$L,P,M_0,d,\zeta,\sigma$ and $\rho_*=M_0^{1-L}$ are now fixed using only
$k$.  Evaluate the explicit minimum above with these numbers to fix $a_0$,
and put
\[
 b_0=\frac{2a_0}{\zeta},\qquad c_k=\frac{a_0}{d}\in(0,1).
\]
The stated range for $c_k$ follows from $0<a_0<\sigma<1$ and $d\ge1$.

Now let $h\in\Z[x]$ be an intersective polynomial of degree $k$ with
positive leading coefficient, and fix its chosen roots as in
Section~\ref{sec:local}.  Apply Lemma~\ref{local:normalization} to $h$.
It supplies positive integers
$D_0,M$, an integer $r$ with $-D_0<r\le0$, and an intersective polynomial
\[
 g(x)=\frac{h(r+D_0x)}{M}\in\Z[x]
\]
of degree $k$ with positive leading coefficient whose chosen roots satisfy
Condition~\eqref{local:unit-condition}.  Moreover,
\[
 Mg(m)=h(r+D_0m),\qquad r+D_0m\ge1\quad(m\in\N_+).
\]

Apply the preceding auxiliary family construction and estimates with $g$ and
its chosen roots as the fixed polynomial data, using the numerical choices
already fixed for degree $k$.  Proposition~\ref{induction:decay} then gives
a constant $C_0\ge1$, which may depend on the fixed normalized polynomial,
its chosen roots, and the fixed structural choices, for which
\eqref{eq:33} holds with the previously chosen $a_0,b_0$.
At $l=1$ the auxiliary polynomial is $g_1=g$, since $r_1=0$ and
$\lambda(1)=1$.  Thus for every positive integer $n$ and every
$B\subseteq[n]$ with $(B-B)\cap g(\N_+)\subseteq\{0\}$,
Corollary~\ref{lift:density} and Proposition~\ref{induction:decay},
both applied to the framework for $g$, give
\[
 \frac{|B|}{n}
 \le Y_1(n,B)^{1/d}
 \le C_0^{1/d}n^{-c_k}.
\]

Now let $N\ge M$ and let $A\subseteq[N]$ satisfy the avoidance condition
of the theorem.  For each $j\in[M]$, set
\[
 A_j=\left\{x\in\left[\left\lceil\frac NM\right\rceil\right]:
            j+M(x-1)\in A\right\}.
\]
These sets give the progression coordinates of a partition of $A$.
If a nonzero difference in $A_j-A_j$ were $g(m)$ for some $m\in\N_+$,
the corresponding difference in $A-A$ would be
$Mg(m)=h(r+D_0m)$, nonzero and with $r+D_0m\ge1$, contradicting avoidance.
The bound just obtained therefore applies to each $A_j$.  Since
$\lceil N/M\rceil\le2N/M$ for $N\ge M$ and $1-c_k>0$, it gives
\[
 |A|=\sum_{j=1}^M|A_j|
 \le M C_0^{1/d}\left\lceil\frac NM\right\rceil^{1-c_k}
 \le 2^{1-c_k}M^{c_k}C_0^{1/d}N^{1-c_k}.
\]
Taking $C_h=2^{1-c_k}M^{c_k}C_0^{1/d}\ge1$ and $N_h=M$ proves the theorem.
The constants $C_h,N_h$ may depend on every coefficient of $h$, while
$c_k$ depends only on $k$.
\end{proof}
